\section{Матрица перехода от базиса к базису. Ее невырожденность. Теорема о преобразовании координат}

\begin{definition}
Пусть в линейном пространстве $V$ заданы два базиса: старый $E = \{e_1, e_2, \dots, e_n\}$ и новый $E' = \{e'_1, e'_2, \dots, e'_n\}$. 
Каждый вектор нового базиса линейно выражается через векторы старого базиса:
\[
\begin{cases}
e'_1 = t_{11}e_1 + t_{21}e_2 + \dots + t_{n1}e_n \\
e'_2 = t_{12}e_1 + t_{22}e_2 + \dots + t_{n2}e_n \\
\vdots \\
e'_n = t_{1n}e_1 + t_{2n}e_2 + \dots + t_{nn}e_n
\end{cases}
\]
Матрица $T$, составленная из коэффициентов этих разложений по столбцам, называется \textbf{матрицей перехода} от базиса $E$ к базису $E'$:
\[
T = \begin{pmatrix}
t_{11} & t_{12} & \dots & t_{1n} \\
t_{21} & t_{22} & \dots & t_{2n} \\
\vdots & \vdots & \ddots & \vdots \\
t_{n1} & t_{n2} & \dots & t_{nn}
\end{pmatrix}.
\]
В блочном виде это записывается как $E' = E \cdot T$.
\end{definition}

\begin{theorem}
Матрица перехода $T$ от одного базиса к другому всегда является невырожденной, то есть $\det T \neq 0$.
\end{theorem}

\begin{proof}
Столбцы матрицы $T$ представляют собой координаты векторов нового базиса $e'_1, \dots, e'_n$ в старом базисе $E$. 
Поскольку $E'$ является базисом, его векторы линейно независимы. 
Известно, что векторы линейно независимы тогда и только тогда, когда линейно независимы их координатные столбцы в любом базисе. 
Матрица, столбцы которой линейно независимы, имеет полный ранг ($\operatorname{rang} T = n$), а следовательно, ее определитель отличен от нуля ($\det T \neq 0$).
\hfill$\blacksquare$
\end{proof}

\begin{theorem}[О преобразовании координат]
Пусть вектор $x \in V$ имеет в старом базисе $E$ столбец координат $X$, а в новом базисе $E'$ --- столбец координат $X'$. Если $T$ --- матрица перехода от $E$ к $E'$, то координаты связаны соотношением:
\[
X = T X'.
\]
\end{theorem}

\begin{proof}
Запишем разложение вектора $x$ по обоим базисам в матричной форме:
\[
x = E \cdot X \quad \text{и} \quad x = E' \cdot X'.
\]
По определению матрицы перехода, $E' = E \cdot T$. Подставим это выражение во второе равенство:
\[
x = (E \cdot T) \cdot X' = E \cdot (T X').
\]
Таким образом, мы получили два разложения одного и того же вектора $x$ по одному и тому же базису $E$:
\[
E \cdot X = E \cdot (T X').
\]
В силу теоремы о единственности координат вектора в данном базисе, столбцы координат должны совпадать. Следовательно:
\[
X = T X'.
\]
\hfill$\blacksquare$
\end{proof}

\begin{remark}
Поскольку $\det T \neq 0$, матрица $T$ обратима. Поэтому формулу преобразования координат можно записать и в обратном виде: $X' = T^{-1} X$.
\end{remark}
